\chapter{Metodologias}

Como primeiro passo, para conseguir viabilizar construir um \textit{software}
com qualidade, é necessário adotar metodologias que fundamentam a ideia e o
processo. De tal sorte que, ajudem a levantar requisitos, priorizar
funcionalidades, quebrar o escopo em entregáveis que possam ser validados, e
desenvolver o aplicativo/sistema de forma sustentável. Esse capítulo se
debruça, tanto na abordagem da metodologia de pesquisa do projeto, quanto
também detalha o método de desenvolvimento de \textit{software} implementado.

\section{Metodologia de Desenvolvimento de \textit{software}}

Existem diversos métodos de desenvolvimento, cada um com seus respectivos prós
e contras. Todos tentam, de sua forma, responder à necessidade de como dar
forma ao processo de desenvolvimento, para que ele flua e assegure que os
requisitos levantados sejam atendidos, que os membros envolvidos no processo,
tanto de desenvolvimento em si, quanto dos atores que o demandam tenham seus
anseios respondidos de forma minimamente transparente e satisfatória.

Além disso, a escolha de uma boa metodologia, adequada para as necessidades do
sistema, permite que os prazos de entrega sejam melhor controlados e mais
precisos \cite{pressman:2021}, minimizando a ocorrência de falhas e gestão do
tempo.

As ideias fundamentais que norteiam esse projeto são baseadas em dois conceitos
chaves da engenharia de \textit{software}: desenvolvimento incremental e
desenvolvimento de \textit{software} ágil. O IEEE define que a engenharia de
\textit{software} é:

\begin{quote}
  Engenharia de \textit{software}: A aplicação de uma abordagem sistemática,
  disciplinada e quantificável no desenvolvimento, na operação e na manutenção de
  \textit{software}; isto é, a aplicação de engenharia ao \textit{software}.
  \cite[p.~67, tradução nossa]{ieee:1990}
\end{quote}

Portanto, a engenharia de \textit{software} propõe estruturar a criação de
\textit{software} de forma mais organizada e eficiente. Pressman vai dar um
destaque para os modelos de desenvolvimento de \textit{software} sequenciais e
incrementais que moldam boa parte das práticas de engenharia de
\textit{software} nas últimas décadas.

O modelo em cascata (\textit{waterfall}) é um modelo sequencial, dos mais
tradicionais e antigos a existirem, cujo qual necessita que cada etapa --
levantamento de requisitos, projeto, implementação, testes -- estejam
concluídas antes da próxima etapa começar. Embora forneça uma estrutura rígida
e lógica, é frequentemente criticado pela baixa adaptabilidade em relação às
mudanças durante o processo de desenvolvimento.

Em outro âmbito, o modelo incremental, trabalha com a premissa de que um
desenvolvimento, feito em pequenas porções, agrupam-se para totalizar o projeto
em si. Trabalho a concepção de incrementos funcionais, permitindo uma maior
flexibilidade com as mudanças no processo ao todo. Também ajuda a iterar
rápidamente em cima do desenvolvimento do sistemas, criando um produto mínimo
viável, que possa preencher a proposta, usando o mínimo de tempo, reduzindo as
funcionalidades ao seu estado basal. Essa metodologia também ajuda a reduzir
custos e riscos, já que necessita que as funcionalidades estejam em seu estado
mínimo, porém funcionais.

Outros modelos tradicionais citados por Pressman incluem o modelo de
prototipação e o modelo em espiral, o primeiro trabalhando na ideia de
elaboração de protótipos pequenos que validem a hipótese, e depois passe para a
implementação final; e o modelo em espiral une a natureza iterativa da
prototipação com os aspectos sistemáticos do modelo em cascata, sendo que acaba
por ser um modelo geralmente usado para desenvolvimento de aplicações em larga
escala, visto que oferece mais segurança ao processo de desenvolvimento. Apesar
da influência histórica, tais modelos acabam por serem um tanto quanto
burocráticos e em um projeto com um escopo de equipe reduzida e onde necessita
um nível de flexibilidade maior, como este, portanto, uma adoção de uma
metodologia mais flexível se faz necessária.

Pressman, comenta que, por meados de 2001, um grupo de desenvolvedores
deflagrou o "Manifesto Ágil", um documento desenhando uma nova metodologia de
desenvolvimento, destacando a satisfação dos patrocinadores do sistema,
entregas mais rápidas, flexibilidade durante o desenvolvimento do
\textit{software} de forma a evitar que os desenvolvedores ficassem empacados e
reduzissem rigidez. Os principais destaques da abordagem ágil é a proposta de
que a mesma não é escrita em pedra, ou seja, permite flexibilidade no como
aplicá-la. Nesse cenário, existem algumas abordagens sobre ela, como o
\textit{Kanban} e o \textit{Scrum}, que dão um pouco mais de forma a ideia
apresentada com o desenvolvimento ágil o que casa relativamente bem com o
modelo de desenvolvimento incremental.

Com base nesse panorama, o presente trabalho adota uma abordagem baseada na
filosofia ágil, estruturada por meio da utilização de um quadro de Kanban,
adaptado à realidade de um desenvolvimento individual. Diferentemente de
modelos tradicionais que envolvem um maior formalismo, o \textit{Kanban}
permite visualizar e gerenciar o fluxo de trabalho de forma contínua, com foco
na entrega incremental de funcionalidades. Por ser um método não prescritivo, o
\textit{Kanban} oferece liberdade para que as tarefas possam ser priorizadas,
ajustadas ou redefinidas conforme a evolução do desenvolvimento, característica
especialmente útil quando há apenas um desenvolvedor responsável por todas as
etapas do processo. Vale ressaltar também, que a metodologia adotada acaba por
ajudar a manter o fluxo de desenvolvimento contínuo sem sobrecarga, já que
limita a quantidade de tarefas em progresso, gerando um movimento de maior
produtividade e entregas.

A escolha pelo \textit{Kanban}, portanto, oferece um equilíbrio sustentável
entre estrutura e flexibilidade, permitindo que o desenvolvimento avance de
maneira organizada, e sem os entraves burocráticos dos modelos mais rígidos.
Essa abordagem permite que o projeto se beneficie dos princípios ágeis, tais
como a entrega contínua de valor, adaptação a mudanças e foco na simplicidade,
ao mesmo tempo que mantém o controle sobre o progresso e a qualidade do produto
final.

\section{Metodologia de Pesquisa}

Entende-se por pesquisa científica, a investigação de um fenômeno, de forma a
encontrar a solução para algum problema \cite{Coelho:2020}. Portanto, a
metodologia científica se torna o conjunto de procedimentos para esta
investigação, cujo qual perpassa a coleta e análise de dados. Neste capítulo,
abordaremos as metodologias adotadas e suas respectivas justificativas, para
além dos dados encontrados para tal.

\subsection{Tipo de Pesquisa}

Segundo \cite{Lakatos:2022}, a pesquisa científica pode ser classificada
quanto aos seus objetivos como exploratória, descritiva ou explicativa. Este
trabalho caracteriza-se principalmente como uma pesquisa aplicada de natureza
exploratória, uma vez que busca investigar práticas, limitações e necessidades
relacionadas a ferramentas digitais para acompanhamento emocional, com o
objetivo de desenvolver uma solução tecnológica prática.

Conforme Jung (2009), a pesquisa aplicada visa gerar conhecimento para
aplicação direta em problemas específicos, com interesse prático e utilitário,
sendo comum em projetos de desenvolvimento de \textit{software}, especialmente
nas áreas de Sistemas de Informação e Ciência da Computação, nas quais o
produto final visa resolver uma demanda concreta do público-alvo.

\subsection{Instrumento de coleta de dados}

Como instrumento de coleta de dados, optou-se pela aplicação de um questionário
estruturado, elaborado com o intuito de compreender hábitos, preferências e
resistências de usuários em relação ao uso de aplicativos voltados para o
acompanhamento emocional e práticas de autorreflexão.

A escolha do questionário se justifica pela sua abrangência, agilidade na
coleta de dados e facilidade de análise quantitativa e qualitativa. Além disso,
como a amostra é composta por usuários potenciais, a aplicação do questionário
de forma digital (via um formulário online) permite alcançar um público mais
amplo com baixo custo e maior comodidade, como também reforça Jung (2009).

O questionário, disponível no Apêndice B, foi bem recebido, e obteve por volta
de 52 respostas válidas. Como a divulgação foi de forma aberta, existia uma
pergunta inicial com uma função chaveadora, a fim de verificar se os
respondentes se encaixavam no perfil de potenciais usuários. Isso se justifica
pela natureza sensível do tema, que envolve saúde mental e experiências
subjetivas.

Na pergunta ``Você já teve algum interesse ou experiência com o monitoramento do
seu humor ou estado emocional?'', apenas 8\% afirmaram que não tem e não teriam
interesse em realizar esse tipo de acompanhamento. Esse dado corrobora outro
resultado relevante: 53\% dos participantes afirmaram já ter, em algum momento
da vida, realizado algum tipo de monitoramento de humor.

Em outros horizontes, mais da metade dos participantes expressou interesse em
acompanhar como o humor varia ao longo do tempo, muitas vezes com sugestões
explícitas como: “Seria legal ver um gráfico semanal ou mensal do meu humor”.
Tal observação reforça a receptividade dos participantes em relação à ideia de
detecção de padrões emocionais, sugerindo que esse tipo de funcionalidade
representa um diferencial positivo frente aos aplicativos existentes.

É notável que os participantes demonstram uma preocupação recorrente com a
privacidade dos dados coletados. Aproximadamente 67\% dos participantes
indicaram desconforto com a exposição de suas informações em plataformas que
não explicam de forma clara como os dados são usados ou armazenados. Essa
percepção está alinhada com a literatura sobre saúde digital. Por exemplo,
\cite{Iwaya:22}, argumentam que a ausência de mecanismos transparentes de
consentimento, controle e anonimização pode minar a confiança dos usuários em
aplicativos de bem-estar e monitoramento contínuo de saúde, especialmente em
contextos sensíveis como saúde mental.

Outro ponto importante revelado pelo levantamento é o desejo por um tom mais
humano, subjetivo e acolhedor. Cerca de 52\% dos participantes relataram
incômodo com ferramentas que adotam uma linguagem genérica, motivacional ou
automatizada demais. Isso reforça a necessidade de uma interface mais empática,
que valorize tanto a escrita livre quanto a análise assistida. Também foram
mencionadas sugestões de funcionalidades consideradas excepcionais, como
inferência automática de ciclos de humor (66\% demonstraram interesse),
personalização da forma de visualização das informações (como gráficos ou
comparativos), e integração com dados contextuais ou ambientais, reforçando a
ideia de um sistema que vá além do simples diário, oferecendo inteligência
emocional útil, mas sem perder a sensibilidade que o tema exige.

\subsection{Análise de Similares}

% TODO: reescrever, arrumar citação
A metodologia trata-se de uma avaliação sobre os sistemas relacionados ao tema
proposto, observando e levantando desafios e possibilidades, validação de
\textit{softwares} e a satisfação dos usuários (Martins, 2018). Seguindo assim,
foram pesquisados sistemas similares encontrados na internet com o objetivo de
realizar uma revisão crítica-reflexiva, analisando os aplicativos que oferecem
algum tipo de diário digital de humor, ou permite algum acompanhamento
reflexivo de variações de emoções.

Para realizar esse processo, de forma a simular também um usuário, foi
pesquisado na Google Play Store \footnote{Google Play Store — Disponível em:
https://play.google.com}, por \textit{mood tracker}, retornando centenas de
resultados. Nesse quesito, estabeleceu os seguintes critérios objetivos: (1)
pode ser um aplicativo em português ou inglês; (2) o aplicativo tem que
apresentar algum tipo de gráfico das entradas feitas (3) tem que possuir
entrada textual.

\subsubsection{Baseline}

O Baseline é um aplicativo de código aberto voltado ao acompanhamento do
bem-estar emocional por meio de uma abordagem minimalista. A ferramenta permite
o registro diário de humor utilizando entradas textuais ou gravações de áudio,
além de disponibilizar notificações personalizadas e revisões periódicas dos
registros realizados.

\begin{figure}[H]
\centering
\caption{Aplicativo Baseline}
\label{fig:screenshot-baseline}
\includegraphics[width=0.40\textwidth,height=0.40\textheight,keepaspectratio]{screenshot-baseline.jpeg}
\source{O autor}
\end{figure}

Entre os aspectos observados, destaca-se a natureza open-source da aplicação,
característica que favorece a transparência e a possibilidade de adaptação por
terceiros. Sua interface apresenta baixo nível de complexidade e reduz a
quantidade de interações necessárias para o registro das informações.
Entretanto, nota-se que o sistema não oferece visualizações temporais avançadas
nem mecanismos de personalização mais aprofundados. Além disso, o
acompanhamento emocional é realizado principalmente por meio de uma pontuação
de humor, sem identificar ciclos emocionais ou relacionar registros a gatilhos
e intervenções específicas.

\subsubsection{DailyBeans}

O DailyBeans é um dos aplicativos de monitoramento de humor mais populares
disponíveis atualmente, possuindo como principal característica a simplicidade
de uso. Nota-se que  possui uma linguagem bem lúdica e informal, por exemplo, o
sistema utiliza o conceito de ``feijões'' representando o estado de humor do
usuário. Destaca-se que o aplicativo permite selecionar categorias de
sentimentos e adicionar atividades para preencher melhor o registro do estado
emocional do usuário para além de um "Feliz", "Triste", etc. Existe também uma
bastante vasta opções de categorias, cada qual com uma ilustração que
representa o que aquilo representa. O aplicativo também apresenta "Temas
sazonais", permitindo que o usuário customize cores, icones, e demais
informações. Um elemento interssante é que as informações adicionadas também
podem ser \textit{offlines}, até uma certa quantidade de entradas/registros,
necessitando de uma conta para poder continuar salvando os dados informados.

\begin{figure}[H]
\centering
\caption{Aplicativo DailyBeans}
\label{fig:screenshot-dailybeans}
\includegraphics[width=0.40\textwidth,height=0.40\textheight,keepaspectratio]{screenshot-dailybeans.jpeg}
\source{O autor}
\end{figure}

Entre suas funcionalidades, além do registro de humor e categorização de
atividades e sentimentos, o aplicativo conta com gráficos e históricos. Uma
interface bastante coesa e visualmente agradável. Um outro elemento é que o
aplicativo permite relacionar indicadores de sono com bastante precisão, como
poder indicar o momento que foi dormir, momento que acordou, anotações, etc. O
aplicativo permite adicionar pelo menos uma foto no plano gratuito, além de
indicar elementos como ``Qual o clima do dia'', se o usuário esteve com
``amigos'', ``familiares''.

Observa-se que o aplicativo possui código fechado e restringe parte de suas
funcionalidades mais avançadas à versão paga ou compras na Loja interna do
aplicativo para obter por exemplo, personalização de icones, adicionar mais
fotos. Além disso, o foco em registros rápidos reduz a flexibilidade para
inserção de reflexões textuais mais detalhadas, e uma exibição excessante de
anúncios atrapalha a utilização.

\subsubsection{DailyMood}

DailyMood é um aplicativo voltado ao registro cotidiano do estado emocional,
oferecendo uma solução enxuta para acompanhamento de humor. Sua proposta
baseia-se em uma interação mínima, permitindo que o usuário selecione
rapidamente o sentimento predominante do dia, sem a necessidade de descrições
extensas. Os registros são organizados em um calendário, o que facilita a
leitura histórica e a identificação de padrões emocionais ao longo do tempo.

\begin{figure}[H]
\centering
\caption{Aplicativo DailyMood}
\label{fig:screenshot-dailymood}
\includegraphics[width=0.40\textwidth,height=0.55\textheight,keepaspectratio]{screenshot-dailymood.jpeg}
\source{O autor}
\end{figure}

Apesar dos recursos disponíveis, observa-se que é uma solução bastante simples.
Não existe outros recursos além de informar o humor predominante e uma
anotação. Um fator interessante é que não necessita de contas, rodando
completamente offline no dispositivo do usuário. Adicionalmente, apresenta
limitações relacionadas à exportação dos dados registrados e à personalização
das informações armazenadas.

\subsubsection{Reflectly}

O Reflectly é um aplicativo de diário pessoal que utiliza recursos de
inteligência artificial para promover práticas de autorreflexão e
acompanhamento emocional. Sua proposta enfatiza a experiência do usuário e a
construção de narrativas pessoais a partir dos registros realizados.

Entre suas funcionalidades estão entradas textuais livres, perguntas guiadas
diárias, análises emocionais automatizadas e personalização visual da
interface. O sistema apresenta uma experiência moderna e intuitiva,
incentivando o uso contínuo por meio de elementos de engajamento.
Um ponto interessante é o fato do aplicativo separar o ``humor'' dos
``sentimentos'', permitindo que os usuários avaliem ambos de forma
independente.

Como pontos observados, destaca-se o fato de ser uma solução proprietária e
dependente de coleta significativa de dados para alimentar seus mecanismos de
análise. Outro ponto negativo é o fato de que o aplicativo apresenta
elementos como estatísticas e histórico, porém necessita de uma conta
paga para acessar esses recursos.

Além disso, sua abordagem está mais voltada ao diário pessoal do que
ao acompanhamento sistemático de indicadores de saúde mental.

\begin{figure}[H]
\centering
\caption{Aplicativo Reflectly}
\label{fig:screenshot-reflectly}
\includegraphics[width=0.40\textwidth,height=0.55\textheight,keepaspectratio]{screenshot-reflectly.jpeg}
\source{O autor}
\end{figure}

\subsubsection{eMoods}

O \textit{eMoods} é um aplicativo de registro de humor especificamente para
pessoas com TAB. O aplicativo tem uma interface bem compacta, bastante
poluída visualmente, em estilo formulário de entrada de dados.

O aplicativo consegue registrar dados de sono, além de perguntas como
ansiedade, estado depressivo, e outras coisas como ``você praticou higiene
pessoal?'', ``você tomou medicamentos?'', ``você se alimentou bem?'' e afins.
Um ponto positivo é que o aplicativo também permite exportar os dados em
formato CSV e gerar relatórios em formato PDF.

Entretanto, a aplicação apresenta limitações relacionadas à experiência do
usuário, possuindo interface considerada desatualizada quando comparada a
soluções contemporâneas. Também não oferece suporte para registros textuais
extensos ou entradas por voz, além de ser muito confuso de usar.

\begin{figure}[H]
\centering
\caption{Aplicativo eMoods}
\label{fig:screenshot-emoods}
\includegraphics[width=0.40\textwidth,height=0.55\textheight,keepaspectratio]{screenshot-emoods.jpeg}
\source{O autor}
\end{figure}

\FloatBarrier

\subsection{Análise comparativa}

\begin{table}[H]
    \centering
    \caption{Comparativo de funcionalidades entre aplicativos analisados}
    \label{tab:comparativo_apps}
    \begin{tabular}{lccccc}
        \hline
        Funcionalidade &
        Baseline &
        DailyBeans &
        DailyMood &
        Reflectly &
        eMoods \\

        \hline
        Registro de humor & Sim & Sim & Sim & Sim & Sim \\
        Registro de sono & Não & Não & Não & Não & Sim \\
        Texto livre & Sim & Parcial & Sim & Sim & Sim \\
        Gráficos & Não & Sim & Sim & Sim & Sim \\
        Gatilhos/intervenções & Não & Parcial & Não & Não & Parcial \\
        Exportação de dados & Sim & Parcial & Parcial & Parcial & Sim \\
        Código aberto & Sim & Não & Não & Não & Não \\
        \hline
    \end{tabular}
    \source{Elaborado pelo autor (2026)}
\end{table}

Com esses dados em mente, podemos compreender alguns elementos fundamentais:
poucos sistemas analisados são \textit{open-source}, o que indica que, em sua
grande maioria, são sistemas comerciais. A maior parte dos sistemas também
oferecem algum aspecto de pagamento, bloqueando funcionalidades importantes
atrás de algum tipo de modelo de subscrição. Existe também o fato de
personalização e de exportação de dados, que poucos ou nenhum sistema oferece,
o que faz com que o sistema não tenha portabilidade dos dados inseridos,
criando assim, um \textit{walled garden}, ou um jardim fechado -- uma técnica de
fazer com que o usuário não consiga migrar entre sistemas, prendendo-o a usar
aquela aplicação já que seus dados estão lá.

\subsection{Relação entre os dados coletados e o desenvolvimento}

A análise dos sistemas similares, quando considerada em conjunto com os dados
coletados por meio do questionário, revela lacunas e oportunidades relevantes
para o desenvolvimento de uma nova solução. Observa-se, por exemplo, que a
maior parte das ferramentas disponíveis no mercado segue modelos comerciais,
frequentemente com funcionalidades essenciais bloqueadas por sistemas de
pagamento ou assinaturas. Além disso, é comum que tais sistemas sejam fechados,
sem disponibilização de código-fonte ou possibilidade de exportação dos dados
inseridos, o que configura uma prática de \textit{walled garden}, dificultando
a migração entre plataformas e restringindo a autonomia dos usuários sobre suas
próprias informações.

Esses aspectos, quando contrastados com o interesse expressivo dos respondentes
por visualizar padrões emocionais ao longo do tempo e por maior personalização,
sugerem direções claras para o desenvolvimento de um sistema que não apenas
atenda a essas demandas, mas também adote princípios de transparência, abertura
e controle por parte do usuário. A seguir, serão discutidas as decisões de
projeto que nortearam a construção da ferramenta proposta, baseadas tanto nas
limitações identificadas nos sistemas existentes quanto nas necessidades e
expectativas manifestadas pelo público-alvo da pesquisa.
